\documentclass[11pt,a4paper]{article}
\usepackage[utf8]{inputenc}
\usepackage[T1]{fontenc}
\usepackage{geometry}
\usepackage{hyperref}
\usepackage{enumitem}
\usepackage{parskip}
\geometry{margin=1in}

\title{\textbf{FlexQL}\\Design and Implementation Report}
\author{Submission for Roll No. 25CS60R65}
\date{April 7, 2026}

\begin{document}

\maketitle

\section{Project Summary}

I implemented FlexQL as a compact SQL-like database server in C++17. My goal was to keep the system small enough to explain end to end while still covering the major pieces of a database driver: parsing, execution, networking, storage, durability, indexing, caching, testing, and benchmarking.

\section{Supported Features}

The implementation currently supports:

\begin{itemize}[leftmargin=1.5em]
    \item \texttt{CREATE TABLE} and \texttt{CREATE TABLE IF NOT EXISTS}
    \item single-row and batched \texttt{INSERT}
    \item \texttt{SELECT *} and projected \texttt{SELECT}
    \item one \texttt{WHERE} condition with \texttt{=}, \texttt{!=}, \texttt{<}, \texttt{>}, \texttt{<=}, and \texttt{>=}
    \item \texttt{INNER JOIN}
    \item single-column \texttt{ORDER BY}
    \item \texttt{DELETE FROM}
    \item row expiry through \texttt{EXPIRES <unix\_ms>}
\end{itemize}

\section{Architecture}

The runtime is split into a small set of modules:

\begin{itemize}[leftmargin=1.5em]
    \item a TCP server that accepts multiple concurrent clients
    \item a client API and simple REPL
    \item a handwritten parser for the supported SQL subset
    \item an executor that performs scans, indexed lookup, joins, and ordering
    \item a storage layer that writes table rows to disk
    \item a durable write-ahead log used for recovery
    \item a bounded primary-key hash index
    \item an LRU cache for repeated \texttt{SELECT} statements
\end{itemize}

The diagrams shipped in the \texttt{diagrams/} directory summarize the same flow from request entry to storage and replay.

\section{Storage and Durability}

I store table rows in disk-backed row files and keep schema metadata separately. For tables with a primary key, I keep a bounded resident hash index and fall back to a disk-backed primary-key hash structure once the table grows beyond the resident limit.

Durability follows a simple WAL-first rule:

\begin{enumerate}[leftmargin=1.5em]
    \item validate the statement,
    \item append a prepared log record,
    \item flush the log,
    \item apply the change locally,
    \item append a commit marker,
    \item flush again,
    \item return success.
\end{enumerate}

On restart, FlexQL replays only committed records and ignores incomplete tails. That keeps crash recovery behavior straightforward and predictable.

\section{Concurrency}

The server uses a worker pool and shared locking around the catalog, tables, cache, and indexes. Reads can proceed in parallel, while mutations are intentionally serialized through the executor's mutation path. I chose the simpler locking model because correctness and recovery were more important than maximum write concurrency for this project.

\section{Testing}

I kept two local test binaries:

\begin{itemize}[leftmargin=1.5em]
    \item \texttt{./bin/flexql-test}
    \item \texttt{./bin/flexql-test-all}
\end{itemize}

The tests cover parser behavior, executor logic, joins, ordering, cache invalidation, TTL handling, durable-log replay, and concurrency smoke checks.

\section{Benchmarking}

I used the integrated benchmark binary for the final benchmark run:

\begin{verbatim}
./bin/flexql-benchmark 10000000
\end{verbatim}

The submission result recorded on April 7, 2026 was:

\begin{itemize}[leftmargin=1.5em]
    \item rows inserted: \texttt{10000000}
    \item insert elapsed: \texttt{9209 ms}
    \item insert throughput: \texttt{1085894 rows/sec}
    \item multi-client full-table \texttt{SELECT}: \texttt{193780 ms}
    \item multi-client primary-key \texttt{WHERE}: \texttt{26198 ms}
    \item unit tests: \texttt{23/23 passed}
\end{itemize}

I saved the full console log in \texttt{benchmark.out}.

\section{Limitations and Future Work}

I deliberately kept the implementation focused. The current system does not include transactions, secondary indexes, aggregate queries, multi-condition boolean predicates, or a cost-based optimizer. If I continue the project, the next improvements would be richer \texttt{WHERE} support, better planning for joins and ordering, and more advanced storage compaction.

\section{Repository Layout}

\begin{verbatim}
flexql/
|-- bin/
|-- include/
|-- src/
|-- tests/
|-- benchmarks/
|-- scripts/
|-- diagrams/
|-- README.md
|-- DESIGN.md
|-- design_doc.tex
`-- benchmark.out
\end{verbatim}

\end{document}
